\documentclass[11pt]{article}
\usepackage[T1]{fontenc}
\usepackage[utf8]{inputenc}
\usepackage{lmodern,amsmath,amssymb,amsthm}
\usepackage[margin=1in]{geometry}
\usepackage[colorlinks=true,allcolors=blue]{hyperref}
\usepackage{microtype}
\newtheorem{theorem}{Theorem}
\newtheorem{lemma}{Lemma}
\newcommand{\Z}{\mathbb Z}
\newcommand{\one}{\mathbf 1}
\newcommand{\HALT}{\mathsf{HALTING}}
\title{A modular-storage proof of NP-completeness for\protect\\ chip-firing halting on simple digraphs}
\author{Alec Kriebel\\\small Independent researcher\\\small ORCID: \href{https://orcid.org/0009-0001-9320-500X}{0009-0001-9320-500X}}
\date{Version 1.0, 16 October 2026}
\begin{document}
\maketitle
\begin{abstract}
We give a self-contained proof that ordinary conservative chip-firing halting
is NP-complete on finite loopless simple strongly connected digraphs.
Hardness holds with polynomial numerical total initial chips and a specified
vertex whose deletion leaves an acyclic graph of depth at most one. A separate
construction forbids antiparallel arcs. The proof combines the primed-toppler
mechanism of Holroyd--Levine--Winkler with the classical prime/Boolean/clause
encoding of Stockmeyer--Meyer, also reconstructed by Gramlich. We prove the
physical conservative realization, its legal checkpoint dynamics and the
arbitrary-schedule implication in full; NP membership is credited to
Farrell--Levine and reproved with an explicit certificate bound. This is an
attributed proof of the historical simple-digraph question. Exact full-result
and realization priority and preexisting worldwide openness remain unknown
after a bounded audit; no firstness is claimed. AI tools were used extensively.
The paper is unrefereed and has not received external human peer review.
\end{abstract}

\section{Question, conventions and attribution}
An input is a finite loopless digraph $G=(V,E)$ with $N=|V|\ge2$, at most one
arc per ordered pair, strong connectivity, and a configuration
$x\in\Z_{\ge0}^{V}$ encoded in binary. A vertex $v$ of outdegree $d_v$ may
fire when $x_v\ge d_v$: it loses $d_v$ chips and sends one along each outgoing
arc. All vertices obey this rule; there is no dissipative sink. A configuration
is stable if $x_v<d_v$ for every $v$. The question $\HALT$ asks whether some
legal finite firing word reaches a stable configuration. Lemma~\ref{least}
below supplies the usual schedule independence, without a fairness assumption.
An \emph{oriented} graph additionally has no antiparallel pair.

Farrell--Levine~\cite[\S3.1, Table~1]{FL} explicitly posed the simple-digraph
NP-completeness question while proving membership in NP and hardness for
binary directed multigraphs. Their historical question is the target here;
an archived open-problem label is not evidence of current worldwide openness.
We establish the following exact statement.

\begin{theorem}\label{main}
$\HALT$ is NP-complete under deterministic polynomial-time many-one
reductions on the input class above. Hardness persists when a specified root
$r$ has the property that $G-r$ contains no directed path of two arcs and is
acyclic, and when the numerical total initial chips is polynomial in the
3-SAT input length. The same assertions hold on the oriented subclass.
\end{theorem}

The ingredients require precise credit. The literal multiplicities in
Lemma~\ref{bank} specialize Holroyd--Levine--Winkler's parallel primed-toppler
construction~\cite[Lemma~4.1]{HLW}. The Boolean residues and local clause
products are the classical Stockmeyer--Meyer encoding~\cite[Theorem~6.1]{SM},
as reconstructed by Gramlich alone~\cite[Fact~3]{Gramlich}, presented here in
complement form. We claim neither encoding as new. Root-last stabilization
and conserved-mass comparison have antecedents in Farrell's thesis
\cite[\S\S1.1,4.1]{Farrell} and Farrell--Levine~\cite[\S2]{FL}.
Least action and tree periods are prior mathematics; our proof of membership
specializes~\cite[\S3]{FL}. Even exponentially long simple-graph games with
polynomial initial chips are prior~\cite{Eriksson}. The argument below supplies
a complete conservative graph realization and SAT equivalence without
asserting priority for that realization or its oriented version.

\section{Least action and a polynomial certificate}
Put $\Delta_{v,v}=d_v$ and $\Delta_{v,u}=-1$ for $u\to v$, with other entries
zero. A word with firing-count vector $g$ changes $x$ to $x-\Delta g$.

\begin{lemma}[Classical least action]\label{least}
If $f\in\Z_{\ge0}^{V}$ has stable $z=x-\Delta f$, every legal word has
firing counts at most $f$ coordinatewise. In particular a nonnegative stable
$z$ is a sound certificate of halting, even if $f$ is not given as a legal word.
\end{lemma}
\begin{proof}
At a first firing that would exceed $f_v$, the current count at $v$ is $f_v$
and all other counts are at most their coordinates of $f$. Nonnegative
incoming arc counts put the chips at $v$ at most
$x_v-d_vf_v+\sum_{u\to v}f_u=z_v<d_v$, contradicting legality.
Thus a legal word has length at most $\sum_v f_v$. Repeatedly firing any legal
vertex must stop at a stable configuration. Conversely, an infinite legal
word rules out every nonnegative stabilizing vector. Early voluntary stopping
at an active state is not halting.
\end{proof}

For completeness we give the polynomial-bit certificate argument, credited to
Farrell--Levine's method. Strong connectivity and the directed matrix-tree
theorem give a strictly positive primitive integer vector $\pi\in\ker\Delta$:
its entries are the numbers of spanning in-arborescences toward each vertex,
divided by their common gcd. At each nonroot tree vertex there are at most
$N-1$ outgoing choices, so
\begin{equation}\label{period}
1\le\pi_v\le(N-1)^{N-1}.
\end{equation}
If $x$ halts, choose its actual terminating firing vector $f$ and final stable
$z\ge0$. Some $w$ has $f_w<\pi_w$: otherwise $f-\pi\ge0$ has the same stable
$z$, and Lemma~\ref{least} applied to the actual word gives
$f\le f-\pi$, a contradiction. For every arc $u\to v$, chip balance yields
\[
 f_u\le\sum_{a\to v}f_a=d_v f_v+z_v-x_v
 \le d_v(f_v+1)-1,
 \qquad f_u+1\le d_v(f_v+1).
\]
Follow a simple directed path \emph{from $u$ to $w$}, of length at most $N-1$.
Since $d_v\le N-1$ and $f_w+1\le\pi_w$, we obtain
\begin{equation}\label{bound}
 f_u+1\le (N-1)^{N-1}\pi_w\le(N-1)^{2N-2}.
\end{equation}
Each coordinate therefore has $O(N\log N)$ bits. For $N=2$ the only graph is
the directed two-cycle; its only nonnegative stable input, by conservation,
is zero, with $f=0$, consistent with~\eqref{bound}. A verifier checks
$f\ge0$ and $0\le x-\Delta f<d$ using binary arithmetic in polynomial time.
This proves membership in NP for both graph classes and arbitrary input chip
magnitudes, without claiming polynomial simulation time.

The additive one above is essential. For arcs $0\to1,1\to0,1\to2,2\to0$,
input $(1,0,0)$ halts after firing $0$ once. The inequality $f_0\le d_1f_1$
would be false. The intermediate source inequality on PDF page~10
of~\cite{FL} omits this slack; the self-contained argument here retains it.
This observation is about that line, not a rejection of the prior membership
theorem or a certification of every other source proof.

\section{An attributed modular bank}
\begin{lemma}\label{bank}
Let $q\ge2$, $\varnothing\ne A\subseteq\Z/q\Z$, and let $h$ be its periodic
indicator on $\Z$. For $0\le a<q$ put
\[
 c_a=1+h(-a-1)-h(-a),\qquad
 S_A(k)=\sum_{a=0}^{q-1}c_a((a+k)\bmod q).
\]
Here residues lie in $\{0,\ldots,q-1\}$. There are $q$ counters in total:
$c_a\in\{0,1,2\}$ and $\sum_a c_a=q$. For every integer $k$,
\begin{equation}\label{storage}
 S_A(k)=K_A+q h(k),\qquad K_A=\frac{q(q-1)}2-|A|.
\end{equation}
Thus its attainable maximum is $B_A=K_A+q$, attained exactly on $A$.
\end{lemma}
\begin{proof}
The multiplicity claims follow from $h\in\{0,1\}$ and telescoping. Advancing
$k$ adds one to each of the $q$ counters, except that a wrapping counter loses
$q$. The wrapping offset is $a=-k-1\pmod q$, whence
\[
 S_A(k+1)-S_A(k)=q-qc_{-k-1}=q(h(k+1)-h(k)).
\]
So $S_A-qh$ is constant. Over one period each counter has average $(q-1)/2$;
the average of $qh$ is $|A|$, proving~\eqref{storage}. Nonemptiness ensures
that the displayed maximum is attained. Full $A$ gives constant storage.
Even if $K_A<0$ in a modulus-two boundary case, all actual residue counts and
storage remain nonnegative. Empty $A$ would not attain $K_A+q$ and is excluded.
\end{proof}

This formula is an exact earlier-mechanism specialization, not a new toppler
lemma. In~\cite[Lemma~4.1]{HLW}, take
$F(k)=k+h(0)-h(k)$ for $k\ge0$. Then $F(0)=0$,
$F(k+q)=F(k)+q$, and its increments are $0,1,2$. The lemma uses
$F(i)-F(i-1)=1+h(i-1)-h(i)$ topplers primed at $a=q-i$; these are precisely
$c_a$. Their cumulative topplings equal $F(k)$. A physical degree-$q$ counter
emits $q$ chips per toppling, so conservation gives
$S_A(k)=S_A(0)+qk-qF(k)$, in agreement with~\eqref{storage}.
A one-output abstract toppler or splitter does not by itself establish an
ordinary conservative chip-firing implementation; the conservative-gate
distinction is explicit in~\cite[\S8.3]{HLW}. We now give that implementation.

\section{The classical Boolean/CRT predicates}
Use an explicitly encoded 3-CNF formula of length $L$, with clauses of length
at most three. Reject malformed encodings, remove duplicate literals, delete
tautological clauses and relabel the variable identifiers actually occurring.
An empty clause is a fixed NO case; no clauses is a fixed YES case. Binary
names or declarations of unused variables are never expanded. Otherwise
there are $n\ge1$ variables, every clause uses $1\le t\le3$ distinct variables,
and the numbers of variables and clauses are $O(L)$.

Choose the first $n$ primes starting at $3$, denoted $p_1,\ldots,p_n$.
For each variable use a bank with modulus $p_i$ and allowed set $\{0,1\}$.
A clause with variable set $I$ uses modulus $q_I=\prod_{i\in I}p_i$ and the
CRT representatives of its satisfying local Boolean assignments. It has
$2^{|I|}-1\le7$ allowed residues, a nonempty proper set. They can be obtained
by inspecting at most eight assignments. An integer $k$ satisfies all bank
predicates exactly when its variable residues describe a satisfying assignment.
Conversely CRT lifts every satisfying assignment to some $k\ge0$ modulo
$M=\prod_i p_i$. This proves both directions without finding an assignment.

For exact attribution, Stockmeyer--Meyer~\cite[Theorem~6.1, pp.~7--8]{SM}
use variable-invalid residues and one failed Boolean residue per clause, with
the same prime and three-prime products; see also~\cite[Fact~3]{Gramlich}.
If $V_i$ means $k\bmod p_i\notin\{0,1\}$ and $F_C$ is the unique falsifying
Boolean residue for clause $C$, their bad language is
\[
 \bigcup_i V_i\ \cup\ \bigcup_C F_C.
\]
Our complete allowed-predicate intersection is exactly its complement. A
local clause complement also contains non-Boolean residues, so it is not
just $F_C$; those additional residues are already in $\bigcup_i V_i$.
This supplies the exact union/complement qualification. Neither this CRT
encoding nor avoidance of expansion of $M$ is claimed as new.

Only individual primes and products of at most three primes are enumerated.
For an explicit polynomial bound, writing $p_j^*$ for the $j$th prime including
$2$, Dusart~\cite[Eq.~(4.2)]{Dusart} gives
$p_j^*\le e j\log j$ for $j\ge2$. Thus
\[
 U:=\max_i p_i=p_{n+1}^*\le3(n+1)^2.
\]
A sieve up to this bound is polynomial; the $j=2$ endpoint also satisfies it.
Every local modulus is at most $U^3$. The global product $M$ supplies only a
correctness witness: the constructor never expands it as a graph or simulates
that many root firings.

\section{A simple conservative realization}
Index the variable and clause banks by $b=1,\ldots,H$, and write
\[
 Q=\sum_b q_b,\quad P=\max_b q_b,\quad
 K_b=\frac{q_b(q_b-1)}2-|A_b|,\quad
 B=\sum_b(K_b+q_b),\quad S(k)=\sum_b S_{A_b}(k).
\]
There are $Q$ \emph{distinct} counters: for each bank and offset $a$, create
$c_a$ separate vertices initially holding $a$ chips. Equal offsets do not
identify vertices. Add a root $r$ and $P-1$ returns $t_1,\ldots,t_{P-1}$.
The complete arc list is
\begin{align*}
 &r\to v &&\text{for every counter }v,\\
 &v\to r,\quad v\to t_j\ (1\le j<q_b)
       &&\text{for a counter of bank }b,\\
 &t_j\to r &&(1\le j<P).
\end{align*}
Hence $d_r=Q$, $d_v=q_b$ and $d_{t_j}=1$. The arc types are disjoint, every
ordered pair occurs at most once, and there are no loops. Root/counter
antiparallel pairs are permitted in this first construction. A modulus-$P$
bank has $P$ actual counters, any of which reaches every return. The root
reaches every counter, and every other vertex reaches the root; thus the
graph is strongly connected. Deleting $r$ leaves only counter-to-return arcs,
an acyclic graph with no two-arc path. The exact counts are
\begin{equation}\label{counts}
 N=Q+P,\qquad |E|=Q+\sum_b q_b^2+P-1.
\end{equation}

Initialize returns at zero and the root at
\begin{equation}\label{rootmass}
 R=Q+B-1-S(0)\ge Q-1\ge0.
\end{equation}
The inequality holds because each bank is at most its own maximum. The total
number of chips is exactly $C=Q+B-1$, and every firing conserves it.
The returns are ordinary degree-one vertices that return their chips to $r$;
they never absorb or delete them.

\section{Legal checkpoints and the SAT equivalence}
Analyze one legal order only for the proof: at a checkpoint all nonroot
vertices are stable and the returns empty. If $r$ is legal, fire it once,
then process all counters, then all returns, before considering $r$ again.
A stable counter gains exactly one chip from that root firing and has no
other incoming arc; it fires at most once, precisely when it reaches $q_b$.
After counters have been processed, each return legally fires once per chip
it holds. These finite phases restore every emitted chip to the root, and
no counter is reactivated before another root firing. Delaying a legal root
while these other legal vertices fire is permitted in ordinary asynchronous
chip-firing. No sink rule has been substituted.

Inductively, at checkpoint $k$, after $k$ root firings, each offset-$a$ counter
holds $(a+k)\bmod q_b$, all returns are zero, and the root holds
\[
 C-S(k)=Q+B-1-S(k).
\]
Using integral storage and $S(k)\le B$ gives the exact equivalences
\begin{equation}\label{checkpoint}
 r\text{ stable}\ \Longleftrightarrow\ S(k)>B-1
 \ \Longleftrightarrow\ S(k)=B
 \ \Longleftrightarrow\ \text{every bank accepts }k.
\end{equation}
A bank misses its maximum by exactly $q_b$ when its predicate is false.
These positive deficits cannot cancel, proving the last equivalence.

If the formula is satisfiable, CRT provides a finite successful $k\ge0$.
The canonical legal process either stops earlier, already proving halting,
or reaches that checkpoint and is stable there. If it is unsatisfiable, each
checkpoint permits a root firing followed by finite legal nonroot phases.
Their concatenation is an infinite legal word. Lemma~\ref{least} rules out
every terminating legal word, regardless of schedule and without fairness.
We have therefore proved
\[
 \text{3-CNF is satisfiable}\quad\Longleftrightarrow\quad
 \text{the constructed ordinary configuration halts}.
\]
The canonical order and possibly enormous stopping time belong to the proof,
not to the polynomial reduction algorithm.

\section{Bounds, totality and the oriented construction}
We have $H=O(L)$, $q_b\le U^3$, $P\le U^3$, and $Q\le HP$.
Equation~\eqref{counts} gives $N=O(HP)$ and $|E|=O(HP^2)$; an explicit
adjacency matrix of $N^2$ bits is polynomial as well. Offset scans and CRT
calculations on at most three primes take polynomial time and bits.
Since $B_b= q_b(q_b+1)/2-|A_b|$,
\[
 0\le C=Q+B-1\le Q+\sum_b\frac{q_b(q_b+1)}2.
\]
The \emph{numerical} total chips and every coordinate are polynomial in $L$,
so unary chip encoding also retains hardness. There is no expansion of a
binary exponential arc multiplicity. As a consistency check, the prior
stable-capacity bound $C\le |E|-N$ for halting inputs~\cite[p.~2]{FL} does
not trivialize the construction: here for every nontrivial input
\[
 (|E|-N)-C=\sum_b\left(\frac{q_b(q_b-3)}2+|A_b|\right)\ge0,
\]
because $q_b\ge3$. This bound is not being used as a SAT decision rule.

For malformed formulas or empty clauses use a directed two-cycle with
configuration $(1,0)$, a NO instance with a circulating chip. For the empty
conjunction use the same graph with $(0,0)$, a stable YES instance. These
cases make the reduction total on encoded strings and preserve all main
graph restrictions, including root deletion. For nontrivial formulas the
construction above applies.

For the oriented version use $P$ returns, omit every counter-to-root arc,
and instead give a modulus-$q_b$ counter the $q_b$ arcs
$v\to t_1,\ldots,v\to t_{q_b}$. Keep $r\to v$ and all $t_j\to r$.
Every arc follows the three-part order root--counter--return--root, so no
opposite pair, loop or repeated ordered pair occurs. Every counter reaches
$r$ through a return, and a largest-bank counter reaches all $P$ returns;
strong connectivity follows. Root deletion again leaves only counter-to-return
arcs. Its exact counts are
\[
 N'=1+Q+P=N+1,\qquad |E'|=Q+\sum_b q_b^2+P=|E|+1.
\]
All outdegrees, counter offsets and root initial mass remain the same.
Each ready counter now emits its $q_b$ chips to $q_b$ distinct returns.
Legally emptying those returns sends exactly $q_b$ chips back to the root,
so the finite-phase proof, residue induction, conservation and
\eqref{checkpoint} all hold with no missing direct output. Thus both SAT
implications and all numerical and bit bounds remain valid. For preprocessing
use the directed three-cycle $r\to s\to t\to r$, with $(0,0,0)$ for YES
and $(1,0,0)$ for NO. Its deletion graph has one arc and its one-chip play
is periodic. This proves the entire oriented case, including trivial inputs.
Together with membership, Theorem~\ref{main} follows.

\section{Priority limits and verification}
The bounded source audit found exact antecedents for the bank priming and CRT
encoding, which are explicitly credited above. It did not establish a source
statement or proved corollary containing the entire restricted theorem.
That is not evidence of worldwide novelty. Priority of the full result,
the conservative simple realization, the oriented realization and any
remaining method contribution remains unknown. No claim that the question
was globally open immediately before this work is made.

Several nearby interfaces must be kept separate. Generalized rotor
reachability~\cite{ACE} has branching arc states and a specified target;
ordinary chip-firing has fixed outputs and asks for any stable endpoint.
Bond--Levine's sandpilization~\cite[Definition~6.5, Corollary~6.6]{BL} preserves
halting on \emph{all} inputs, not automatically one reset-state input.
Hardness for containing weighted-net classes does not automatically restrict
to one consuming transition per place, degree thresholds, unit outputs,
conservation and a supplied marking. Esparza's sole-author survey~\cite{Esparza}
does not remove those hypotheses. The full Mayr--Weihmann construction and
complete Hujsa--Devillers weighted-net proofs were unavailable to this audit;
their possible additional implications remain unknown. Optimization over
added chip distributions, infinite backgrounds and binary arc expansion are
also different decision interfaces. These statements record exact transfer
gaps, rather than exclude every possible earlier derivation.

The accompanying portable package contains this source, bibliographic hashes
and read scopes, frozen authored history, new standard-library exact finite
controls and a deterministic archive rebuilder. It distributes no third-party
PDFs or copied source bodies. New diagnostics check bank identities and the
two exact antecedent substitutions, legal conservative main/oriented graphs,
small arbitrary schedules, certificate slack, and deliberately damaged
variants. They supplement the general proof, do not certify the literature,
and are not replays of historical omitted programs. Source editions and
selective read scopes are distinguished explicitly in the source ledger;
the actual journal and arXiv editions of~\cite{ACE} have different bytes.

AI tools were used extensively in solving, checking, adversarial review,
source comparison, programming and writing. The author is responsible for
the claims. This manuscript is unrefereed, has not received external human
peer review or journal acceptance, and is not formally certified in a proof
assistant. No individual was contacted. The paper and authored verification
materials are licensed under Creative Commons Attribution 4.0 International.

\begin{thebibliography}{99}
\bibitem{FL} M. Farrell and L. Levine, \emph{CoEulerian graphs},
Proc. Amer. Math. Soc. 144 (2016), 2847--2860.
\href{https://doi.org/10.1090/proc/12952}{doi:10.1090/proc/12952}.
Inspected edition: \href{https://arxiv.org/abs/1502.04690v3}{arXiv:1502.04690v3}.
\bibitem{HLW} A. E. Holroyd, L. Levine and P. Winkler,
\emph{Abelian logic gates}, inspected
\href{https://arxiv.org/abs/1511.00422v2}{arXiv:1511.00422v2} (2018),
Lemma~4.1, p.~15, and \S8.3, p.~32.
\bibitem{SM} L. J. Stockmeyer and A. R. Meyer,
\emph{Word Problems Requiring Exponential Time: Preliminary Report},
Proceedings of the Fifth Annual ACM Symposium on Theory of Computing (1973),
Theorem~6.1 and proof, pp.~7--8 of the retained preliminary report.
\bibitem{Gramlich} G. Gramlich, \emph{Probabilistic and Nondeterministic Unary
Automata}, MFCS 2003, LNCS 2747, 460--469. Fact~3, PDF pp.~5--6 of the
\href{https://ggramlich.github.io/Publications/unaryMFCS03.pdf}{author-hosted manuscript}.
\bibitem{Farrell} M. Farrell, \emph{The halting problem for chip-firing on
finite directed graphs}, Cornell senior thesis, 2 December 2014,
\href{https://pi.math.cornell.edu/files/Research/SeniorTheses/farrellMatthewThesis.pdf}{author thesis},
\S\S1.1,4.1. L. Levine is identified as thesis advisor, not coauthor.
\bibitem{Eriksson} K. Eriksson, \emph{No Polynomial Bound for the Chip Firing
Game on Directed Graphs}, Proc. Amer. Math. Soc. 112(4) (1991), 1203--1205.
\bibitem{Dusart} P. Dusart, \emph{Estimates of Some Functions Over Primes
without R.H.}, \href{https://arxiv.org/abs/1002.0442v1}{arXiv:1002.0442v1}
(2010), Eq.~(4.2), PDF p.~4.
\bibitem{ACE} D. Auger, P. Coucheney and K. R. Etse,
\emph{Linear Extensions of Rotor-Routing in Directed Graphs: Reachability
Problems}, Electron. J. Combin. 33(2) (2026), P2.5,
\href{https://doi.org/10.37236/13580}{doi:10.37236/13580}, Theorem~21 and
Lemmas~22--23, actual journal PDF pp.~39--43.
\bibitem{BL} B. Bond and L. Levine, \emph{Abelian networks II. Halting on all
inputs}, inspected \href{https://lionellevine.github.io/abelian-networks-II-halting.pdf}{author-hosted manuscript},
Definition~6.5 and Corollary~6.6, PDF pp.~17--18.
\bibitem{Esparza} J. Esparza, \emph{Decidability and complexity of Petri net
problems---an introduction}, inspected
\href{https://lsv.ens-paris-saclay.fr/~demri/Esparza98.pdf}{sole-author survey},
selected definitions and class interfaces, PDF pp.~31--32.
\end{thebibliography}
\end{document}
